% =============================================================================
\section{Security Implications}
\label{sec:implications}
% =============================================================================

\subsection{Adversarial Use: DPRI as a Reconnaissance Tool}

An adversary with access to multiple candidate datasets---for example,
a researcher evaluating which of several publicly released models to
attack---can compute DPRI in minutes without any model training and
prioritize high-risk targets.

We simulate this scenario.
An adversary has a budget of $B$ MIA queries (each query = training and
attacking one model on one dataset).
Under random dataset selection, the expected AUC return is the
population mean.
Under DPRI-guided selection (attack highest-DPRI first), the adversary
reaches the top-$k$ risk datasets in expectation with
\TODO{X}$\times$ fewer queries (Figure~\ref{fig:adv_query}).

\begin{figure}[t]
  \centering
  \includegraphics[width=0.85\linewidth]{../results/regression/adv_query_efficiency.pdf}
  \caption{Adversarial query efficiency: DPRI-guided vs.\ random dataset
           selection. Shaded region = 95\% CI over 1,000 random orderings.}
  \label{fig:adv_query}
\end{figure}

\paragraph{Mitigation.}
Data holders can compute DPRI before releasing a dataset or a model
trained on it.
A high DPRI score should trigger additional safeguards: stronger DP
guarantees, data minimization, or access controls on the trained model.

\subsection{DP Budget Recommendation}

Based on the empirical relationship between DPRI and DP efficacy
(Finding~3), we propose a simple DPRI-adaptive $\varepsilon$ rule:

\begin{equation}
  \varepsilon_{\mathrm{rec}}(D) =
  \begin{cases}
    \varepsilon_0 & \text{if } \mathrm{DPRI}(D) < 0.35 \\
    \varepsilon_0 / 3 & \text{if } 0.35 \leq \mathrm{DPRI}(D) < 0.65 \\
    \varepsilon_0 / 10 & \text{if } \mathrm{DPRI}(D) \geq 0.65
  \end{cases}
  \label{eq:eps_rec}
\end{equation}

where $\varepsilon_0$ is the practitioner's nominal budget.
This rule is calibrated so that all datasets achieve approximately
equal post-DP MIA AUC, regardless of their native structure.
The thresholds are derived from our empirical calibration on
Table~\ref{tab:dp_calibration} and should be recalibrated as more
datasets are added to the benchmark.

\subsection{Policy Implication}

Current regulatory frameworks for machine learning privacy (e.g., GDPR
\S{}22, HIPAA de-identification) treat privacy as a property of the
algorithm, not the data.
Our findings imply that this is insufficient: two organizations applying
identical DP-SGD with identical $\varepsilon$ achieve systematically
different levels of membership inference protection depending on the
native structure of their data.

We recommend that data governance frameworks adopt DPRI (or a
standardized equivalent) as a required pre-training assessment,
analogous to environmental impact assessments required before
construction projects.
The DPRI tool we release makes this practical at near-zero computational
cost.

\subsection{Extensions}

\paragraph{Federated learning.}
In federated settings, each client's local dataset has its own DPRI.
Clients with high DPRI should contribute more noise in the local DP
step, or be excluded from the federation if their risk exceeds a
threshold.

\paragraph{Generative models.}
DPRI is defined on the training data distribution and is
model-agnostic; it applies equally to GANs, diffusion models, and LLMs.
We leave empirical validation on generative model training data as
future work.
